% LSPV-4 REPORT::PATHSTRAT1::2026-10-04T07:49:00-07:00
\documentclass[11pt]{article}
\usepackage[margin=.85in]{geometry}
\usepackage{amsmath,amssymb}
\begin{document}
\section*{LSPV-4 PATHSTRAT1: realized-path versus strategy-output dichotomy}
\textbf{STATUS: OPEN.} P versus NP remains open. The theorem below is \textbf{PROVED/EXACT} under its explicit hypotheses and is a red-team classification, not a SAT lower bound.

\section*{CORPUS REVIEWED}
The current corpus was surveyed through repository search and direct readback. Governance read: AGENTS.md, CODEX.md, and the pnp-research skill. Current root main.tex and verified recovery MAIN0658 were read. Causal sources read or searched include TASK-004, TASK-005, canonical CCSPC, CONTGAME1, LOCALSEM1, ROBUSTGAP1, COLLAPSEGATE1, and DEFGATE1. Structural/acquisition state was checked through root main and the active TASK-001 references. The newest LSPV-3 IDENTITYGATE1 body is absent from GitHub and was not reconstructed from its fallback summary. The predecessor failure is nonblocking.

\section*{Target}
DEFGATE1 shows that no solver-independent continuation problem $W_{\rm exact}$ is yet source-defined. This run asks what happens to the weakest natural definitions.

Let a polynomial-horizon CCSPC instance have realized transcript
\[
\tau=(p_0,\Phi_0,T_0,B_0,Y_0,\ldots,p_r),
\]
where each realized stage transition and audit response is computable in deterministic polynomial time with polynomially many calls to SAT on polynomial-size formulas. Call a staticized task \emph{path-faithful} if its required output/correctness concerns only the one challenge sequence encoded or realized in the instance. Call it \emph{strategy-strengthened} if correctness requires responses for counterfactual legal future challenges, for example through
\[
\forall B_i\exists Y_i\forall B_{i+1}\exists Y_{i+1}\cdots .
\]

\section*{Theorem PATHSTRAT1}
Under the stage hypothesis above:

\paragraph{(A) Realized-path side.}
Every path-faithful polynomial-horizon staticization whose output can be generated by sequentially executing those stage procedures lies in $\mathrm{FP}^{\mathrm{SAT}}$. Hence, conditionally on $P=NP$, it lies in FP.

\paragraph{(B) Strategy side.}
Canonical realized-path CCSPC correctness does not by itself imply correctness for a strategy-strengthened task. Any implication from arbitrary static SAT correctness or canonical realized-path CCSPC to such an all-futures task requires an additional source/solver-obligation theorem.

\paragraph{Proof.}
For (A), execute the polynomially many realized stages in order. A polynomial number of stages, each using polynomial local work and polynomially many polynomial-size SAT queries, composes to one polynomial-time SAT-oracle transducer. Thus the output is in $\mathrm{FP}^{\mathrm{SAT}}$. If $P=NP$, deterministic polynomial-time SAT queries can be replaced by polynomial-time subroutines, giving FP.

For (B), canonical CCSPC specifies one realized challenge stream and explicitly separates it from the strengthened every-challenge objective. A machine may wait for the actual $B_i$ and compute only its corresponding $Y_i$. Counterfactual branches are absent from that correctness condition. Therefore the universal future obligations require a new theorem; they are not obtained by definition-preserving staticization. \(\square\)

\section*{Why this is not merely DEFGATE1}
DEFGATE1 establishes a specification gap: $W_{\rm exact}$ is not source-defined. PATHSTRAT1 classifies the two weakest repair directions once a definition is proposed. Preserving realized-path semantics leaves a polynomial SAT-oracle composition under the stated stage hypothesis. Adding all-futures strategy semantics can escape that closure only by strengthening the obligation, and the source/solver-obligation arrow remains OPEN.

This yields the hostile fork
\[
\boxed{
\begin{array}{ll}
\text{path-faithful }W_{\rm exact} &\Rightarrow \mathrm{FP}^{\mathrm{SAT}}\text{ under the stage hypothesis},\\[2mm]
\text{strategy-strengthened }W_{\rm exact} &\Rightarrow \text{new obligation theorem required}.
\end{array}}
\]

\section*{Primary literature}
Chandra, Kozen, and Stockmeyer, \emph{Alternation}, JACM 28(1), 1981, prove alternating polynomial time equals deterministic polynomial space. This is a \textbf{KNOWN THEOREM}; it classifies genuine all-futures alternation but does not force static SAT to realize it.

Chew and Clymo, ECCC TR19-157 / CADE 2020, prove that polynomial-time strategy extraction for QRAT would imply $P=PSPACE$; universal expansion is the crucial source of hardness in their declared proof system. This is a \textbf{KNOWN THEOREM} and a model-boundary warning, not an arbitrary-SAT transfer theorem.

Chew, AAAI 2025, formalizes round-based QBF strategy extraction into circuits for specified QBF proof systems and derives proof simulations. Again, the strategy representation is supplied by the proof system rather than forced on arbitrary SAT deciders.

\section*{Hostile controls}
\paragraph{RECOMP and preprocessing.}
On the path-faithful side, RECOMP is exactly compatible with the proof: wait for the realized challenge and recompute its response. Full static preprocessing can only make this easier. On the strategy side, preprocessing does not create a source obligation to answer counterfactual challenges unless those challenges are part of the defined problem.

\paragraph{Compact strategies.}
Explicit strategy-tree size is not a lower bound. If an all-futures strategy has a polynomial-size evaluator, the project must charge that evaluator and test PH/round absorption. If no compact evaluator exists, that requires an independent lower bound in the matching representation.

\paragraph{Wrappers.}
DEAD, COPY, and DELAY can alter trace length and chronology without changing either side of the dichotomy. INTERP can replace the native realization. Therefore implementation-specific trace geometry cannot close the missing obligation theorem.

\paragraph{Semantic and algebraic oracles.}
Giving SAT answers for free preserves the semantic gap between one realized path and all counterfactual futures. Algebraic/low-degree access likewise does not manufacture the missing universal obligation. PATHSTRAT1 therefore relativizes and supplies no nonalgebrizing hinge.

\paragraph{Easy families.}
Horn, 2-SAT, affine/Tseitin, and bounded-treewidth controls fit the same fork. Easy stage evaluation makes the path side easier; quantified strategy semantics, when independently imposed, is a different problem. No formula-side structural hardness is promoted to arbitrary runtime.

\section*{Structural branch audit}
Nothing here weakens or strengthens the SU/vMS/Hadamard program. Exact TASK-001 source transcription, quotient-normalized characteristic-two rank growth, wrapper-invariant source-matched acquisition, representation-invariant no-sharing, and fully charged reconstruction remain OPEN. A structural theorem still requires the arbitrary-algorithm transfer.

\section*{DUPLICATION / CONTRADICTION AUDIT}
PATHSTRAT1 adopts DEFGATE1, LOCALSEM1/ROBUSTGAP1, CONTGAME1, PH/round closure controls, and WRAPINV1 at their verified scopes. It does not reconstruct WDEFGATE1 or IDENTITYGATE1. It sharpens the queue by showing that the most conservative path-preserving definition falls into SAT-oracle composition under explicit stage hypotheses, while the stronger strategy definition reopens the source-obligation gap.

No contradiction with COLLAPSEGATE1: if a separately defined strategy target is PSPACE-hard and a universal polynomial compiler exists, COLLAPSEGATE1 states the resulting class-collapse consequence.

\section*{WHAT CHANGED RELATIVE TO ROOT MAIN}
Root main currently says to define the weakest solver-independent continuation object, attack staticization/RECOMP, prove the directional source relation, then classify it. PATHSTRAT1 resolves the first hostile case split under explicit stage hypotheses: path-preserving definitions are SAT-oracle compositions; all-futures definitions require a new obligation theorem.

\section*{ADOPTED / REJECTED / RETIRED PRIOR COMPONENTS}
Adopted: DEFGATE1, ROBUSTGAP1, COLLAPSEGATE1, CONTGAME1, WRAPINV1, SAME\_REALIZATION discipline, and strict charge separation.

FALSIFIED/RETIRED as generic hinges: claiming that merely naming a path-faithful $W_{\rm exact}$ escapes P-closure under $P=NP$; claiming that universalizing future challenges is semantics-preserving without an obligation theorem; explicit strategy-tree size; freshness alone; trace chronology alone.

\section*{Barrier audit}
PATHSTRAT1 relativizes. It does not cross algebrization. Natural Proofs is not directly triggered because no broad constructive circuit property is proposed. A future strategy-size lower bound would require a fresh Natural-Proofs audit if it becomes a broad efficiently recognizable circuit property.

\section*{Components consistent with $P=NP$}
The path-faithful closure is explicitly consistent with $P=NP$. The semantic distinction between realized paths and all-futures strategies is also consistent with $P=NP$. Current SU/vMS fixed-realization algebra, quotient normalization, provenance closure controls, and reconstruction ledgers remain consistent with $P=NP$ at their declared scopes.

\section*{OPEN arrows}
No load-bearing arrow closes. OPEN: exact source-defined $W_{\rm exact}$; a source/solver-obligation theorem for any strategy-strengthened target; independent membership and hardness classification; universal static-SAT compiler; exact TASK-001 rank growth; wrapper-invariant acquisition; TC4; representation-invariant no-sharing; fully charged reconstruction; nonrelativization; nonalgebrization.

\section*{Single next falsification target}
Write the weakest externally fixed solver-independent $W_{\rm exact}$ that is not merely path-faithful, and state the exact obligation theorem
\[
\forall A\in{\cal A}^{P}_{SAT}\quad
\operatorname{CorrectSAT}(A)\Longrightarrow
\operatorname{Solves}(C_A,W_{\rm exact}).
\]
Attack this implication first with RECOMP and full preprocessing, then compact strategy evaluators, INTERP, maximal SAT-oracle replacement, and algebraic-oracle replacement. Do not classify $W_{\rm exact}$ before its source obligation is explicit.

\section*{Final status}
P versus NP remains \textbf{OPEN}. PATHSTRAT1 is a local classification theorem under explicit stage hypotheses, not a separation.

% END_OF_VERIFIED_REPORT::LSPV-4::PATHSTRAT1
\end{document}
